% Scene Composition Without Grammar Extension:
% Plugin-Based Domain Vocabulary for Spatial Languages
% Target: ACM I3D 2027 (submit Nov '26)
% Backup venue: SIGGRAPH Asia 2027 (submit May '27)
% Format: ACM sigconf
% Status: Draft — threats + USD comparison expanded; measured LOC table pending (2026-04-18)
% Paper ID: P3-S3 (third in Program 3: The Stalk and the Center)
% Subsystem proven: packages/core/src/parser/HoloCompositionParser.ts + packages/plugins/
% Draft: April 2026

\documentclass[sigconf,nonacm]{acmart}

% ── Packages ─────────────────────────────────────────────────────────────────
\usepackage{amsmath}
\let\Bbbk\relax % acmart/newtxmath already defines \Bbbk; prevent amssymb redefinition clash
\usepackage{amssymb}
% Cross-paper macros: \hsplus and \compilefn used in body but not defined
% locally; provide minimally. Other shared macros (\holo, \hashfn,
% \tropplus, \tropmult, \provsemiring) are defined by paper-12's own
% \newcommand block lower in this preamble — don't shadow.
\providecommand{\hsplus}{\texttt{.hsplus}}
\providecommand{\compilefn}{\mathsf{compile}}
\usepackage{booktabs}
\usepackage{xcolor}
\usepackage{listings}
\usepackage{algorithm}
\usepackage{algpseudocode}
\usepackage{tikz}
\usetikzlibrary{arrows.meta,positioning,shapes.geometric,fit,calc,
  decorations.pathreplacing,backgrounds}
\usepackage{pgfplots}
\pgfplotsset{compat=1.18}
\usepackage{subcaption}
\usepackage{balance}
\usepackage{stmaryrd}

% ── Draft annotations ────────────────────────────────────────────────────────
\newcommand{\todo}[1]{\textcolor{red}{\textbf{[TODO: #1]}}}
\newcommand{\stub}[1]{\textcolor{gray}{\textit{[stub: #1]}}}

% Lotus petal data contract — \measuredFrom{<artifact>} marks values that come
% from a runnable artifact in the repo. Renders nothing in the PDF; grep target
% for scripts/build-lotus.mjs (research/2026-04-26_lotus-petal-data-contract.md §3+§4).
\providecommand{\measuredFrom}[1]{}

% ── Convenience macros ──────────────────────────────────────────────────────
\newcommand{\holoscript}{\textsc{HoloScript}}
\newcommand{\holo}{\texttt{.holo}}
\newcommand{\hashfn}{\mathcal{H}}
\newcommand{\tropplus}{\oplus}
\newcommand{\tropmult}{\otimes}
\newcommand{\provsemiring}{\mathcal{S}}
\newcommand{\pluginset}{\mathcal{P}}
\newcommand{\domainbridge}{\textsc{DomainBridge}}

% ── Theorem environments ────────────────────────────────────────────────────
\newtheorem{definition}{Definition}
\newtheorem{property}{Property}
\newtheorem{theorem}{Theorem}
\newtheorem{lemma}{Lemma}
\newtheorem{corollary}{Corollary}

% ════════════════════════════════════════════════════════════════════════════
\title{Scene Composition Without Grammar Extension:\\
Plugin-Based Domain Vocabulary for Spatial Languages}

\subtitle{Provenance-Preserving Domain Extension\\
in the \holoscript{} \holo{} Format}

\author{Joseph Krzywoszyja}
\affiliation{%
  \institution{Infinitus Systems}
  \country{USA}
}
\email{brianonbased@gmail.com}

% ════════════════════════════════════════════════════════════════════════════
\begin{document}

\begin{abstract}
Scene description languages face a domain-extension dilemma: every
new application domain (medical simulation, robotics, film VFX,
structural biology) demands domain-specific vocabulary, but
extending the core grammar to accommodate each domain leads to
language fragmentation and maintenance collapse. We present
\holo{}, a scene composition format in which domain vocabulary
enters as \emph{plugins}, not as core grammar extensions. The core
\holo{} grammar recognizes exactly five structural forms (objects,
templates, spatial groups, traits, metadata) and never changes; domain
plugins register trait categories, object archetypes, and validation
rules through a typed extension interface. We prove a
provenance-preservation theorem: the hash chain from source \holo{}
through plugin-applied domain vocabulary to compiled output is
unbroken. The theorem witnesses a two-implementation agreement: the
forward path (compiler emits the chain as a side effect of producing
the compiled artifact) and the backward path (the verifier walks the
hash decomposition and recomputes each link from the same source +
plugin inputs) are independent implementations of the same provenance
state machine, and the theorem guarantees deterministic replay across
the pair through identical hash outputs at every level of the
core+plugin decomposition --- the same-adapter strict-digest regime
of the W.315 wire-format equivalent primitive, with the Route~2b
per-step canonical projection pattern handling cross-vendor
cases~\cite{paper-0c-twin-test-2026-05-11}. The theorem builds on
database and workflow provenance
models~\cite{buneman2001provenance,prov-dm-w3c}
and provenance-semiring algebra~\cite{green2007provenance}, but
applies them to spatial language extension rather than query
evaluation. Unlike OpenUSD schema plugins, which extend the data
model without compositional trust guarantees, \holo{} keeps the
extension trace in the compiled artifact. We discuss shipped
\holoscript{} plugins (AlphaFold, medical, robotics, film-VFX,
legal-document, Hololand) and a multi-target compile pipeline with
no core grammar edits; baseline (plugin-free) wall-clock compile provenance
overhead is frozen in
\texttt{packages/core/src/compiler/\_\_tests\_\_/paper-targets-overhead-bench.test.ts}
with logs in \texttt{research/benchmark-results-2026-04-17-paper-targets-overhead.md}
(\texttt{.ai-ecosystem}) and \texttt{HoloScript/.bench-logs/}. Plugin-specific
overhead versus non-plugin scenes follows the same W.067 policy as Paper 10:
only targets above the percent-overhead floor receive ratio claims. Phase 3
bench extension (commit \texttt{6d4d1186d}) adds 18 additional CompilerBase
targets as below-floor exclusions with archived logs and frozen summary in
\texttt{research/benchmark-results-2026-04-17-phase3-noise-floor.md}.
This is the third
paper in Program~3 of a 16-paper research
program~\cite{trustbyconstruction2026,capstone-uist2027}.
\end{abstract}

\keywords{scene description, plugin architecture, domain extension,
  provenance, grammar composition, OpenUSD, spatial computing,
  digital twins}

\maketitle

% ────────────────────────────────────────────────────────────────────────────
\section{Introduction}
\label{sec:intro}

Every spatial computing standard faces the domain-extension
dilemma. VRML established a web-oriented 3D scene graph lineage,
and X3D (ISO 19775) carried that lineage into a profile-based
modularity system~\cite{carey1997vrml,brutzman2008x3d}. X3D
partitions the specification into profiles---Interchange,
Interactive, Immersive---but adding a new domain (e.g., medical
simulation) requires proposing a new profile through ISO, a
multi-year standards process.
glTF (Khronos) allows \texttt{KHR\_}/\texttt{EXT\_}
extensions that modify the serialization schema, but these
create vendor-specific forks: a file with
\texttt{EXT\_mesh\_gpu\_instancing} is not a valid glTF file
on implementations that lack the extension. OpenUSD (Pixar/
AOUSD) supports schema plugins via \texttt{plugInfo.json}
registration, but plugin-contributed prims carry no
provenance annotation across the extension
boundary~\cite{aousd-core-spec-1.0}.

In all three systems, adding domain vocabulary either changes
the core language (X3D profiles), fragments the ecosystem
(glTF extensions), or produces provenance-invisible
contributions (USD schema plugins). We present \holo{}, a
scene composition format whose core grammar is
\emph{immutable}---five structural forms, formally specified,
unchangeable across all domain applications. Domain knowledge
enters as a \emph{plugin} that registers vocabulary within
the existing forms. The core parser sees the plugin's
contributions as instances of its five forms; it does not
need to be modified, extended, or recompiled.

\paragraph{Contributions.}
\begin{enumerate}
  \item The \holo{} core grammar: five structural forms, formally
    specified, immutable across all domain applications.
  \item The plugin extension interface: typed registration of trait
    categories, object archetypes, and validation rules. Plugins
    extend vocabulary without touching the parser.
  \item Provenance-preservation theorem: the hash chain spans the
    plugin boundary. A compiled output from a plugin-extended
    \holo{} file traces back to both the core source and the
    plugin's contribution, composed via the provenance semiring.
  \item Case studies on shipped plugins: AlphaFold (structural
    biology), medical simulation, robotics (Isaac Sim), film-VFX
    (volumetrics, GCode), legal-document (programmable law), and
    Hololand (game world). Zero core grammar modifications in each case study.
  \item Comparison with OpenUSD schema plugins: provenance breaks
    at the USD plugin boundary because USD's composition arcs
    (LIVERPS) resolve \emph{value opinions}, not \emph{trust
    chains}.
\end{enumerate}

\paragraph{Novelty.}
To our knowledge, \holo{} is the first scene composition format to
provide formal guarantees that domain extension preserves the
provenance chain end-to-end. X3D, glTF, and OpenUSD all support
some form of domain extension (profiles, vendor extensions, schema
plugins, respectively), but none carry hash-chain identity through
the extension boundary: a prim contributed by a USD schema plugin
is provenance-invisible to the USD composition arc resolver.
\holo{} is also the first scene language to treat the five
structural forms as an \emph{invariant}: the grammar surface cannot
be enlarged by any downstream consumer, guaranteeing a single
authoritative parser across all domain deployments.

% ────────────────────────────────────────────────────────────────────────────
\section{Background and Related Work}
\label{sec:related}

\subsection{Scene Description Languages}
X3D (ISO 19775) defines a profile-based modularity system:
the specification is partitioned into Interchange, Interactive,
and Immersive profiles, each adding features to the prior.
Profiles are spec-level partitions decided by standards
committees, not user-level plugins. The VRML-to-X3D
migration fractured the ecosystem along profile boundaries.
glTF (Khronos) uses a JSON-based asset format with
\texttt{KHR\_} (ratified) and \texttt{EXT\_} (vendor)
extensions. Extensions modify the serialization schema;
implementations that do not support an extension may reject
or silently ignore extension data. Neither approach carries
provenance through the extension boundary.
This missing boundary record matters because provenance systems are
useful precisely when they can answer what changed, where it came
from, and how it was transformed~\cite{w3c-prov,Herschel2017}.

USD's schema plugin system is discussed in
\S\ref{sec:related} above.

\subsection{Language Extension Mechanisms}
Racket's \texttt{\#lang} mechanism allows user-defined
languages within a host language. Haskell's type classes
provide open extension via instance declarations. Elixir's
protocol system enables polymorphic dispatch over data types
without modifying the data type definition. All three
provide extension mechanisms for general-purpose programming
languages. None are designed for scene description, and none
carry provenance annotations through the extension boundary.
Compiler infrastructures such as MLIR show that dialect extension can
scale domain-specific computation when the extension boundary is
explicit~\cite{mlir-lattner2020}; \holo{} adopts the same separation
discipline at the scene-language layer, where the extension unit is a
domain plugin rather than a compiler dialect.
The \holo{} plugin model differs in that the extension
interface is constrained to the five structural forms:
plugins cannot introduce new syntactic categories, only new
vocabulary within existing categories.

% ────────────────────────────────────────────────────────────────────────────
\section{The \holo{} Core Grammar}
\label{sec:grammar}

\subsection{Five Structural Forms}

The \holo{} core grammar defines exactly five structural forms:

\begin{enumerate}
  \item \texttt{object $id$ \{ $\overline{p}$ \}}: A typed
    spatial entity with named properties. Object types
    (\texttt{mesh}, \texttt{light}, \texttt{camera}, etc.)
    are vocabulary, not grammar---new types are registered
    by plugins.
  \item \texttt{template $id$ ($\overline{x}$) \{ $\overline{p}$ \}}:
    A parameterized blueprint that produces objects upon
    instantiation. Templates are first-class values.
  \item \texttt{group $id$ \{ $\overline{e}$ \}}: A spatial
    hierarchy container. Groups define parent-child
    transform relationships.
  \item \texttt{trait $id$ : $category$ \{ $\overline{p}$ \}}:
    A behavioral annotation from the \hsplus{} trait
    system~\cite{paper-11-hsplus-ecoop2027}. Trait categories
    (\texttt{physics}, \texttt{animation}, \texttt{medical},
    etc.) are vocabulary, not grammar.
  \item \texttt{metadata \{ $\overline{k{:}v}$ \}}: Key-value
    annotations for tooling, provenance, and regulatory
    metadata. Metadata blocks participate in the hash chain.
\end{enumerate}

The parser is implemented in
\texttt{HoloCompositionParser.ts} and
recognizes exactly these five forms. Domain-specific keywords
(e.g., \texttt{protein}, \texttt{joint}, \texttt{scalpel})
are resolved at the vocabulary level, not the grammar level.

\subsection{Immutability Invariant}
\begin{property}[Grammar Immutability]
No plugin may add, remove, or modify the five core
structural forms. Plugins may register new vocabularies
\emph{within} existing forms: new object type names, new
trait category names, new validation functions. The parser's
production rules are fixed at compile time and do not depend
on which plugins are loaded. This invariant is enforced by
the parser's architecture: plugin contributions are resolved
\emph{after} parsing, during a separate semantic pass.
\end{property}

% ────────────────────────────────────────────────────────────────────────────
\section{The Plugin Extension Interface}
\label{sec:plugins}

\subsection{Registration Model}

A plugin registers four categories of domain vocabulary via
a typed TypeScript interface:

\begin{enumerate}
  \item \textbf{Object type names.} A set of domain-specific
    type strings (e.g., \texttt{protein}, \texttt{ligand},
    \texttt{scalpel}). The parser accepts these as valid
    \texttt{object} type values without modification.
  \item \textbf{Trait categories.} A set of domain-specific
    trait category names (e.g., \texttt{foldable},
    \texttt{dockable}, \texttt{surgical}). Trait composition
    follows the algebra of~\cite{paper-11-hsplus-ecoop2027}.
  \item \textbf{Validation functions.} Domain-specific
    constraint checkers that run during the semantic pass
    (e.g., verifying that atomic coordinates are within
    physiological ranges for a medical plugin).
  \item \textbf{Compilation hooks.} Target-specific lowering
    extensions that produce domain-aware output (e.g., the
    AlphaFold plugin emits PDB-format metadata alongside
    the glTF structural output).
\end{enumerate}

The parser sees every plugin contribution as an instance of
the five core forms. A \texttt{protein} object is parsed as
an \texttt{object} with type \texttt{protein}; the parser
does not know or care that \texttt{protein} is a plugin
contribution.

\subsection{Provenance at the Plugin Boundary}

When a plugin contributes a trait or object type, its
contribution carries a provenance annotation. The hash of
the compiled output includes the plugin's identity. We model this as
a provenance semiring in the tradition of semiring-weighted
provenance and shortest-distance algebra~\cite{green2007provenance,mohri2002semiring},
with the tropical operator matching the min-plus lineage used by the
compiler core~\cite{mikhalkin2005tropical}:
$$\hashfn(\text{output}) = \hashfn(\text{core source})
  \tropmult \hashfn(\text{plugin}_i)$$
The plugin's identity hash $\hashfn(\text{plugin}_i)$ is
derived from the plugin's registered name, version, and
validation function signatures. An output compiled from a
medical \holo{} file is cryptographically distinguishable
from an output compiled from a robotics \holo{} file---by
provenance, not by naming convention.

\subsection{Provenance-Preservation Theorem}
\begin{theorem}[Plugin Provenance Preservation]
For any core \holo{} source $s$ and plugin set
$\mathcal{P} = \{p_1, \ldots, p_k\}$, the compiled output
$\compilefn(s, \mathcal{P}, T)$ for target $T$ carries a
hash that decomposes as:
$$\hashfn(\compilefn(s, \mathcal{P}, T)) =
  \hashfn(s) \tropmult
  \hashfn(p_1) \tropmult \cdots \tropmult \hashfn(p_k)
  \tropmult \hashfn(T)$$
No plugin can contribute silently; all contributions are
algebraically visible in the hash chain.
\end{theorem}

\begin{proof}[Proof sketch.]
By induction over the plugin application order. Each
plugin $p_i$ contributes vocabulary that maps to one of the
five core forms. The parser evaluates each form according to
the operational semantics of~\cite{paper-10-hs-core-pldi2027},
which threads the provenance hash. The plugin's identity hash
enters as a $\tropmult$ factor at the point where the
plugin's vocabulary is resolved during the semantic pass.
Since $\tropmult$ is associative and commutative for
non-interfering plugins, the decomposition holds regardless
of application order. Two plugins that register the same
object type name produce a conflict detected at registration
time, not at compilation time.
\end{proof}

% ────────────────────────────────────────────────────────────────────────────
\section{Case Studies}
\label{sec:cases}

\subsection{AlphaFold Plugin (Structural Biology)}
Registers protein, ligand, and chain object types;
structural-biology traits (\texttt{foldable},
\texttt{dockable}); and validation functions for atomic
coordinate ranges (C$\alpha$ bond length $\in [1.4, 1.6]$\AA).
Core grammar: untouched. Parser diff: zero lines.
Reference: \texttt{packages/plugins/alphafold-plugin/}.

\subsection{Medical Simulation Plugin}
Registers anatomy object types (\texttt{organ},
\texttt{vessel}, \texttt{tissue}), surgical-procedure traits
(\texttt{incisable}, \texttt{suturable}), and patient-safety
validation rules (instrument reach constraints, hemodynamic
bounds). Core grammar: untouched.
Reference: \texttt{packages/plugins/medical-plugin/}.

\subsection{Robotics Plugin (Isaac Sim)}
Registers robot-arm object types (\texttt{joint},
\texttt{link}, \texttt{end\_effector}), joint-constraint
traits (\texttt{revolute}, \texttt{prismatic}), and Isaac
Lab compatibility validation (Xform ordering, rotateXYZ
conventions). Core grammar: untouched.
Reference: \texttt{packages/plugins/robotics-plugin/}.

\subsection{Film-VFX Plugin}
Registers volumetric object types (\texttt{vdb\_volume},
\texttt{particle\_system}), GCode slicing traits for
3D-printed props, and VFX compositing validation (layer
ordering, alpha premultiplication). Core grammar: untouched.
Reference: \texttt{packages/plugins/film3d-volumetrics-plugin/}.

\subsection{Additional Plugins}
Two further plugins demonstrate the model's generality:
\begin{itemize}
  \item \textbf{Legal-document plugin.} Registers contract
    object types and regulatory-compliance traits for
    programmable law. Enables \holo{} files that describe
    legal instruments with hash-verified provenance.
  \item \textbf{Hololand plugin.} Registers game-world
    object types (\texttt{npc}, \texttt{quest\_trigger},
    \texttt{spawn\_point}) and gameplay traits. Extends
    the VRR compilation target with world-state metadata.
\end{itemize}
Both: zero core grammar modifications.

% ────────────────────────────────────────────────────────────────────────────
\section{Evaluation}
\label{sec:eval}

\subsection{Parser and Compiler Regression Suites}
Plugin extensibility assumes the core grammar never forks. In CI, parser and
compiler behavior are anchored by focused Vitest suites under
\texttt{packages/core/src/parser/} (including \texttt{HoloCompositionParser.test.ts}
and domain-block tests) and \texttt{packages/core/src/compiler/\_\_tests\_\_/},
which replay representative \texttt{.holo}/\texttt{.hsplus} inputs and multi-target compile
pipelines.

\subsection{Grammar Immutability Verification}
Across shipped plugins, we verify that the core \holo{}
parser has zero modifications. The metric is the line-level
diff of \texttt{HoloCompositionParser.ts} between the
plugin-free base build and each plugin-loaded configuration.

\paragraph{Parser immutability check.}
The release process runs a line-level diff of
\texttt{HoloCompositionParser.ts} across shipped plugin
configurations; the expected result is zero parser edits
attributable to plugins (plugins register via the stable API only).

\subsection{Compilation Overhead}
We compile a representative suite of \holo{} scenes (from single-object domain files to large multi-plugin hierarchies)
with and without plugins loaded, across the supported compile
targets. The metric is wall-clock compile-time overhead:
$$\text{overhead}(\%) = \frac{t_{\text{plugin}}
  - t_{\text{base}}}{t_{\text{base}}} \times 100$$

\paragraph{Compile overhead.}
For plugin-free multi-target compile overhead, cite the 2026-04-17 frozen
harness (\texttt{benchmark-results-2026-04-17-paper-targets-overhead.md};
\texttt{paper-targets-overhead-bench.test.ts}): prefer $\Sigma$-run aggregates
on fast code targets and the glTF row for stable positive percent overhead in prose.
On structural grounds, the cost
concentrates in plugin validation hooks rather than registration
(validation walks the AST; registration is one dispatch-table insert).

\subsection{Reproducibility Pipeline and Calibration}
\label{sec:paper12-repro-pipeline}

The Paper~12 calibration path is a two-layer replay. The core provenance
overhead claim is rerun from the HoloScript repository root with
\texttt{pnpm --filter @holoscript/core exec vitest run}
\texttt{src/compiler/\_\_tests\_\_/paper-targets-overhead-bench.test.ts};
the frozen prose summaries are
\texttt{research/benchmark-results-2026-04-17-paper-targets-overhead.md}
and
\texttt{research/benchmark-results-2026-04-17-phase3-noise-floor.md}.
The plugin-loaded scene-suite and OpenUSD comparison are rerun with
\texttt{pnpm --filter @holoscript/comparative-benchmarks exec vitest run}
\texttt{src/\_\_tests\_\_/paper-12-scene-suite-overhead.bench.test.ts}
and
\texttt{src/\_\_tests\_\_/paper-12-structural-biology-comparison.bench.test.ts}.
Those harnesses regenerate
\texttt{HoloScript/.bench-logs/2026-04-27-paper-12-scene-suite-overhead.md}
and
\texttt{HoloScript/.bench-logs/2026-04-27-paper-12-structural-biology-comparison.md}.
Reviewer calibration is therefore mechanical: rerun the two commands, verify
the generated artifacts, and refresh this \LaTeX{} source only when the
artifact values move.

\subsection{Complexity, Cost, and Scalability Model}
\label{sec:paper12-cost-scaling}

For a parsed scene with $N$ core nodes, $P$ loaded plugins, $D$ plugin-authored
domain nodes, and $T$ compile targets, the decoder cost is
$O(N + P + D)$ for parsing plus plugin registration and $O(T(N + D))$ for
target lowering with provenance-chain emission. Memory is $O(N + P + D)$
because plugin identity hashes are stored once and referenced from each
domain contribution. The OpenUSD comparison has a different fixed cost:
schema authoring, \texttt{usdGenSchema}, plugin registration, CMake, install,
environment binding, and host-code invocation before a scene is decoded.

\paragraph{Provenance decoder cost.}
Verification of the provenance-preservation theorem (the hash chain spans the plugin boundary) costs $O(P + N_\text{total})$ under the semiring folding, where $P$ is the number of plugins and $N_\text{total}$ the aggregate plugin-contributed nodes. A P/N sweep (P\in\{1,2,4,8,16\}, N\in\{10,100,1000\}) on the chain-walking decoder is recorded in \texttt{research/paper-12-holo-decoder-cost.md}; the empirical envelope matches the inductive bound derived from the theorem. (Grok Build, task\_1779176007458\_1cr6).

The pinned camera-ready envelope is intentionally modest and repeatable:
the scene-suite matrix covers 5 scenes across 2 target paths, 1--20 objects,
and 0--4 traits per object; the 2026-04-27 artifact reports that the largest
cell parses in $<0.45$\,ms cold and exports plugin-authored USDA in
$<0.08$\,ms. Scale-out is by independent scene shards rather than by mutating
the grammar: batch replay can scale to million-scale operations because each
scene's core parse, plugin registration record, target output, and hash-chain
verification are partitionable artifacts. The core grammar therefore stays
fixed while validation throughput grows with the number of replay workers, not
with new syntax.

\subsection{Provenance Chain Integrity}
For each compiled output from each plugin configuration,
we verify that the hash chain decomposes into core + plugin
contributions as the provenance-preservation theorem
predicts. The decomposition check is:
$$\hashfn(\text{output}) \stackrel{?}{=}
  \hashfn(s) \tropmult
  \hashfn(p_1) \tropmult \cdots \tropmult \hashfn(p_k)
  \tropmult \hashfn(T)$$

\paragraph{Chain integrity.}
The decomposition check must pass for every plugin configuration
on the regression scene suite; any violation is a release blocker.

\subsection{Comparison with OpenUSD Schema Plugins}
We implement the same domain-extension scenario (structural
biology: protein objects with foldability traits) in both
\holo{} and USD. The comparison measures:
\begin{enumerate}
  \item \textbf{Extension effort.} USD requires:
    \texttt{schema.usda} authoring, \texttt{usdGenSchema}
    code generation, \texttt{plugInfo.json} registration,
    and C++ plugin compilation. \holo{} requires: a single
    TypeScript registration call.
  \item \textbf{Provenance visibility.} USD's compiled
    output (a \texttt{.usd} file) does not carry a
    provenance chain through the schema boundary: a prim
    created by the biology plugin is indistinguishable in
    hash terms from a built-in prim. \holo{}'s compiled
    output carries the plugin's identity in its hash chain.
\end{enumerate}

\paragraph{USD comparison.}
On structural grounds, the effort differential is order-of-
magnitude: USD requires schema authoring, \texttt{usdGenSchema} code
generation, \texttt{plugInfo.json} registration, and a C++ compile
round trip (hundreds of lines and multiple toolchain steps per standard
USD workflow), versus a single typed registration call plus trait
definitions on the \holo{} side.

\subsection{Threats to Validity}
\label{sec:threats}

\paragraph{Plugin representativeness.}
The shipped plugins stress diverse domains but do not exhaust every
OpenUSD or glTF deployment pattern; compile-time overhead and provenance
checks could differ on pathological scene sizes or extremely deep template
expansion.

\paragraph{Toolchain drift.}
OpenUSD and compiler stacks evolve on independent cadences. The USD
comparison table must name versions (\texttt{usdGenSchema}, core library)
so LOC and step counts remain reproducible.

\paragraph{Trust boundary.}
Plugins are trusted at registration time; behavior outside that model falls
to Program~1 sandboxing~\cite{paper-4-sandbox-usenix} rather than to the
\holo{} grammar itself.

% ────────────────────────────────────────────────────────────────────────────

\subsection{GPU Benchmark and Ablation Study}
\label{sec:ablation}

\paragraph{Plugin-free baseline.}
Compiling \holo{} scenes with no plugins loaded establishes the
baseline compile time for the five-form grammar alone.
Incrementally loading the AlphaFold, medical, robotics, and
film-VFX plugins in sequence isolates the per-plugin registration
cost. Expected behavior: registration cost scales linearly with
plugin count (one dispatch-table insert per plugin); provenance-chain
depth scales with plugin-contributed node count; and the core parser
line-diff remains zero throughout.

\paragraph{Provenance overhead ablation.}
The 2026-04-17 frozen harness
(\texttt{paper-targets-overhead-bench.test.ts}) ablates the hash
path by compiling with and without \texttt{provenanceHash} enabled
across eleven targets. For the fast code targets (sub-5\,ms base
median), overhead is indistinguishable from timer noise; the glTF
row, at a 210\,ms base median, yields a stable $+$11.57\% (Phase~1)
and $+$16.60\% on $\Sigma$-run --- the most defensible figure for a
positive percent claim under W.067 policy. The RTX~6000 Ada rerun
(2026-04-24; frozen in
\texttt{research/benchmark-results-2026-04-24-core-rtx6000ada.md})
confirms the Phase~1 results within JIT noise bands across hardware.

\paragraph{Full-loop demonstration.}
The reference demonstration pipeline exercises the system
end-to-end: (1)~author a \texttt{.holo} file that imports from the
AlphaFold and robotics plugins simultaneously; (2)~invoke the parser
and verify grammar immutability with a line-level diff;
(3)~compile to WebGPU and USD targets in a single pipeline invocation;
(4)~extract the provenance hash chain and verify the decomposition
theorem holds for each output. This pipeline is fully automated as
the CI smoke-test suite in
\texttt{packages/core/src/compiler/\_\_tests\_\_/}, covering
multi-plugin scenes, multi-target compile, and provenance chain
decomposition, and runs unmodified on any checkout from commit
\texttt{548a42ff}.

\subsection{User Study Protocol (Preregistered)}
\label{sec:user-study}

To evaluate the operational utility and cognitive ergonomics of the
plugin model, we have preregistered a user study ($N=12$) targeting
participants who extend a scene with a new domain vocabulary under
two conditions: (A)~the \holo{} plugin registration interface and
(B)~the OpenUSD schema-plugin workflow. A preliminary pilot study
($N=3$) validated the experimental harness and confirmed the
feasibility of the survey instrument. The full study measures task
completion time for a plugin registration task, error recovery rates
when the participant introduces a domain-vocabulary conflict, and
qualitative cognitive load (NASA-TLX) against the OpenUSD baseline.
The study protocol, task script, IRB waiver, and preregistration
packet are archived in \texttt{research/user-study-kit/} to satisfy
D.011 reproducibility requirements. Camera-ready results replace
this paragraph once the $N=12$ cohort completes.

\section{Discussion}
\label{sec:discussion}

\subsection{Limitations}

\paragraph{Trusted plugins.}
The plugin model assumes plugins are trusted code. A
malicious plugin could register object types that pass
validation but produce semantically incorrect outputs.
Sandboxing plugin execution---preventing malicious plugins
from corrupting the scene graph---is addressed by
Program~1's sandbox paper~\cite{paper-4-sandbox-usenix}.

\paragraph{Language-boundary overhead.}
Plugins written in languages other than TypeScript (e.g.,
Python for scientific computing, Rust for performance-critical
domains) require an FFI bridge. The provenance overhead of
cross-language plugin calls is not yet measured.

\paragraph{Plugin composition conflicts.}
Two plugins that register the same object type name produce a
conflict detected at registration time. The current model
rejects conflicts; a future extension could support
namespace-qualified type names.

\subsection{Future Work}

\begin{itemize}
  \item \textbf{Plugin marketplace.} Domain-specific trait
    categories distributed via x402 micropayments, enabling
    an ecosystem of paid domain plugins with provenance-verified
    contributions.
  \item \textbf{Dynamic plugin loading.} Currently, plugins
    are registered at compile time. Runtime plugin loading
    would enable dynamic domain extension without
    recompilation.
  \item \textbf{Plugin composition algebra.} Extend the
    P3-S2 trait conflict algebra~\cite{paper-11-hsplus-ecoop2027}
    to the plugin level: when two plugins contribute
    conflicting trait categories, resolve the conflict
    algebraically rather than by rejection.
\end{itemize}

\subsection{Remaining Work for Camera-Ready}
\begin{itemize}
  \item \textbf{Plugin-loaded measurements (scene-suite matrix).} The
5-scene suite (tiny / small / medium / large / plugin-heavy; 1--20
objects, 0--4 traits per object) is run on the canonical replay host
across two target paths: HoloScript parser cold/warm and the
\texttt{@holoscript/openusd-plugin} export pipeline. Measured suite
aggregates are HoloScript cold parse mean/median/p95/max =
0.192 / 0.180 / 0.448 / 0.448\,ms; warm parse =
0.130 / 0.108 / 0.280 / 0.280\,ms; OpenUSD plugin export =
0.033 / 0.028 / 0.077 / 0.077\,ms. The largest cell (20 objects, 4
traits) parses in $<$0.45\,ms cold and exports plugin-authored USDA
in $<$0.08\,ms; cold/warm ratio stays within 0.5--0.9 across the
suite, indicating amortized parser memoization is the dominant warm-
path effect rather than per-trait composition cost.
\measuredFrom{HoloScript/packages/comparative-benchmarks/src/__tests__/paper-12-scene-suite-overhead.bench.test.ts}%
\measuredFrom{HoloScript/.bench-logs/2026-04-27-paper-12-scene-suite-overhead.md}
  \item \textbf{Structural-biology extension comparison (USD pin v25.11).}
The structural-biology extension is authored on both sides: HoloScript
ships \texttt{@holoscript/structural-biology-plugin} (3 object types
\texttt{protein}/\texttt{ligand}/\texttt{chain} + 4 annotation traits
\texttt{foldable}/\texttt{helix}/\texttt{sheet}/\texttt{residue\_anchor}
registered through a single typed \texttt{register()} call) and the
OpenUSD comparison authors the equivalent schema-plugin pinned to
upstream tag \texttt{PixarAnimationStudios/OpenUSD~v25.11} (commit
\texttt{363a7c8d}). Effective code LOC: \textbf{HoloScript~=~113}
(\texttt{src/index.ts}~98 + \texttt{package.json}~15) vs
\textbf{OpenUSD~=~323} hand-authored (\texttt{schema.usda}~131 +
\texttt{plugInfo.json}~86 + \texttt{tokens.h}~33 +
\texttt{CMakeLists.txt}~73), giving a USD/Holo authored ratio of
\textbf{2.86$\times$} before the additional 600--1200 LOC that
\texttt{usdGenSchema} codegen produces (\texttt{wrap*.cpp},
\texttt{module.cpp}, per-class \texttt{.h}/\texttt{.cpp} pairs,
\texttt{generatedSchema.usda}). Toolchain steps: \textbf{HoloScript~=~1}
(call \texttt{register(host)}); \textbf{OpenUSD~v25.11~=~8} (author
\texttt{schema.usda}, author \texttt{plugInfo.json}, run
\texttt{usdGenSchema}, author \texttt{CMakeLists.txt}, configure CMake,
build+install, set \texttt{PXR\_PLUGINPATH\_NAME}, author host code).
Provenance visibility: the HoloScript plugin fuses plugin id and
version into the per-residue anchor and per-object chain hash, so a
downstream consumer can recover plugin attribution from the compiled
artifact alone (\texttt{verifyChain} rejects trait tampering
end-to-end on a 12-residue test protein); on the OpenUSD side a
compiled \texttt{.usdc} loses plugin attribution after stage flatten,
because schema composition is not algebraically fused with prim
identity in v25.11.
\measuredFrom{HoloScript/packages/plugins/structural-biology-plugin/src/index.ts}%
\measuredFrom{HoloScript/packages/comparative-benchmarks/usd-comparison/structural-biology/schema.usda}%
\measuredFrom{HoloScript/packages/comparative-benchmarks/src/__tests__/paper-12-structural-biology-comparison.bench.test.ts}%
\measuredFrom{HoloScript/.bench-logs/2026-04-27-paper-12-structural-biology-comparison.md}
\end{itemize}

% ────────────────────────────────────────────────────────────────────────────
\section{Conclusion}
\label{sec:conclusion}

The domain-extension dilemma---change the core language or
fragment the ecosystem---is a false dichotomy. \holo{}
resolves it by separating grammar from vocabulary. Five
immutable structural forms host arbitrary domain vocabularies
via a typed plugin interface. The provenance semiring spans
the plugin boundary: every domain contribution carries an
algebraic annotation that makes it visible in the hash chain.

Shipped domain plugins---structural biology, medical
simulation, robotics, film-VFX, legal documents, game
worlds---illustrate the model across radically different
application domains, without core grammar edits. One verifiable chain.
The third stalk of the lotus.

% ────────────────────────────────────────────────────────────────────────────
% Bibliography migrated 2026-04-24 from inline \begin{thebibliography} block
% to shared research/holoscript.bib per founder /founder ruling Decision 2
% (research/paper-audit-matrix.md §2026-04-24 Refresh — Founder Ruling).
% Pattern documented in docs/paper-bibliography-migration.md.
% Inline block preserved below inside `\iffalse` guard for reviewer audit;
% remove at submission bundle time after holoscript.bib has been validated
% to resolve every \cite key this paper uses.
\bibliographystyle{ACM-Reference-Format}
\bibliography{holoscript}

\iffalse
% === LEGACY INLINE BIBLIOGRAPHY (pre-2026-04-24 migration) ===
% Kept behind \iffalse for audit trail only; bibtex ignores it.
% Remove entirely at submission-bundle time after:
%   pdflatex paper-12-holo-i3d && bibtex paper-12-holo-i3d
% resolves 0 missing keys against research/holoscript.bib.
\begin{thebibliography}{30}

\bibitem{trustbyconstruction2026}
J.~Krzywoszyja.
\newblock Trust by Construction: Provenance-Native Simulation Contracts.
\newblock \emph{IEEE TVCG}, 2026 (submitted).

\bibitem{capstone-uist2027}
J.~Krzywoszyja.
\newblock From Notation to Cognition: A Provenance-Native Architecture.
\newblock In \emph{UIST}, 2027 (forthcoming).

\bibitem{paper-10-hs-core-pldi2027}
J.~Krzywoszyja.
\newblock HoloScript Core: A Contracted Compilation IR for Spatial Computing.
\newblock In \emph{PLDI}, 2027.

\bibitem{paper-11-hsplus-ecoop2027}
J.~Krzywoszyja.
\newblock Reactive Traits Under Contract: Interaction Semantics for Spatial Languages.
\newblock In \emph{ECOOP}, 2027.

\bibitem{paper-4-sandbox-usenix}
J.~Krzywoszyja.
\newblock Sandboxing Spatial Computation.
\newblock In \emph{USENIX Security}, 2026.

\bibitem{paper-13-dumbglass-siggraph2028}
J.~Krzywoszyja.
\newblock Dumb Glass: Rendering as Contracted Synthesis of All Projections.
\newblock In \emph{SIGGRAPH}, 2028 (forthcoming).

\bibitem{green-semiring2007}
T.~J.~Green, G.~Karvounarakis, and V.~Tannen.
\newblock Provenance semirings.
\newblock In \emph{PODS}, 2007.

\bibitem{aousd-core-spec-1.0}
Alliance for OpenUSD (AOUSD).
\newblock OpenUSD Core Specification 1.0.
\newblock Ratified December 17, 2025.
\newblock \url{https://aousd.org/specifications/}

\bibitem{x3d-iso19775}
ISO/IEC 19775.
\newblock Extensible 3D (X3D) architecture.

\bibitem{gltf-khronos}
Khronos Group.
\newblock glTF 2.0 Specification.
\newblock \url{https://registry.khronos.org/glTF/specs/2.0/glTF-2.0.html}

\bibitem{racket-lang}
M.~Flatt and PLT.
\newblock The Racket Language.
\newblock \url{https://racket-lang.org/}

\bibitem{wadler-typeclasses1989}
P.~Wadler and S.~Blott.
\newblock How to make ad-hoc polymorphism less ad hoc.
\newblock In \emph{POPL}, 1989.

\bibitem{elixir-protocols}
J.~Valim et al.
\newblock Elixir protocols and polymorphism.
\newblock \url{https://elixir-lang.org/}, 2024.

\end{thebibliography}
\fi

\end{document}
